\question{What does FLIP-IT compute, and where does each step happen?}
\label{sec:method}
\statusMethod

\begin{shortanswer}
Every computation that touches a patient row happens inside the practice. Each round, a practice
receives the current global coefficient vector, trains on its own rows for two passes, and returns
an updated vector plus the number of examples it trained on. The server multiplies those vectors by
their example counts, sums, and divides. No patient row, and no derived per-patient quantity, is
ever transmitted --- including the feature scaler, which is fitted locally and stays there.
\end{shortanswer}

\subsection{The local pipeline}

A practice's SuperNode holds its own extract and never sends it anywhere. Per round, in order:

\begin{enumerate}
  \item \textbf{Load} the practice's rows. In the pilot this is the practice's own FHIR server or
        its Tomedo export; in the test bed it is a generated CSV. The loader is the only component
        that differs, deliberately, so that everything downstream is identical either way.
  \item \textbf{Split} 80/20 into train and held-out, by a permutation seeded from the global seed
        and the practice index --- reproducible, and stable across rounds so that the held-out rows
        are never trained on.
  \item \textbf{Scale} features with a \texttt{StandardScaler} fitted \emph{on this practice's
        training rows only}. This matters for the disclosure analysis: the scaler encodes practice
        means and variances, and it is \textbf{never transmitted}. Each practice keeps its own.
  \item \textbf{Train} logistic regression, warm-started from the global vector, for two full
        passes, with balanced class weighting for the prevalence imbalance.
  \item \textbf{Reply} with the updated coefficient vector and \texttt{num-examples}.
\end{enumerate}

The model is logistic regression over \NumFeatures{} features plus an intercept: age, six binary
comorbidity flags, and three ``years since diagnosis'' durations. \S\ref{sec:modelchoice} explains
why this model class rather than a stronger one, and why that choice is also the privacy-preferable
one.

\subsection{The round trip}

FLIP-IT runs on Flower \cite{beutel2020flower}, version \FlwrVersion{}. In the deployment runtime the
\textbf{SuperLink} is a message broker: practices (\textbf{SuperNodes}) pull instruction messages and
push reply messages; the ServerApp retrieves the replies and aggregates them. The aggregation rule
is sample-weighted federated averaging \cite{mcmahan2017fedavg}: with $w_i$ the vector returned by
practice $i$ and $n_i$ its example count,
\[
  w_{\text{global}} \;=\; \frac{\sum_i n_i\, w_i}{\sum_i n_i}.
\]
\texttt{num-examples} is exactly the $n_i$ above --- it is transmitted because the average is
weighted by it, and for no other purpose.

Two properties of this arrangement drive the whole of \S\ref{sec:isolation}, and they are worth
stating here rather than as a surprise later: \textbf{aggregation is the last step of a round, not a
property of transmission}, and \textbf{every message carries its sender's identity}. Neither is a
defect; both are simply how the protocol works, and both are what secure aggregation exists to
change.

\subsection{Configuration}

The complete training and aggregation configuration --- the third of the four items the assessment
listed as undocumented --- is a single declarative block in the project's \texttt{pyproject.toml},
reproduced in full in Appendix~\ref{sec:config}. It is not scattered through code, and it is
versioned with the code it configures, so any run's configuration is recoverable from the commit
that produced it.
